\section{Experiments}
\label{sec:experiments}

\subsection{Experimental setup}

\paragraph{Inference modes.}
We evaluate models in two settings:
(1)~\textbf{S3-only}: models receive a ground-truth corrected frame containing
the key and must read it. This tests pure terminal-state extraction and serves
as an upper-bound control slice.
(2)~\textbf{Full pipeline}: models receive the scrambled \texttt{.mp4} video
and natural-language instructions with Bash access, and must independently
parse instructions (S1), execute \texttt{ffmpeg} corrections (S2), locate the
key window, and extract the key (S3).

The distinction is central: strong performance in S3-only does not imply strong
inverse program recovery. It only shows that once the correct terminal state is
provided, the model can decode the target string. Full-pipeline performance
measures whether the system can actually reach that terminal state.

\paragraph{Agentic evaluation.}
For full-pipeline evaluation, VLM agents are given Bash access and must decide
their own video processing strategy (ffmpeg flags, frame sampling, temporal
search). No ground-truth metadata (operations, key, timing) is provided.

\paragraph{Reported metrics.}
We report strict exact match (EM), character-level F1, stage-specific
S1 parse F1, S2 execution success rate, and Pipeline-F1. In the current draft,
the strongest supported evidence comes from S3-only and calibrated extraction
experiments; full-pipeline reporting remains incomplete and is therefore
presented as an explicit gap rather than as a filled leaderboard.

For the video-temporal diagnostic domain, we additionally report \emph{process
partial credit}. This score decomposes a failed answer into operation credit,
frame-localization credit, displayed-fragment similarity, normalized-answer
similarity, transform-derived-answer similarity, final-answer similarity, and
tool-use credit. It is intentionally not a replacement for EM; it is a dense
diagnostic signal used to distinguish uninformative failures from meaningful
near misses.

\subsection{Models}

We benchmark OCR baselines, heuristic pipelines, and frontier multimodal
systems:

\begin{description}
  \item[\textbf{OCR-only}]
  Pure extraction baseline. Uses easyocr on sampled key frames; no language
  understanding, program inference, or tool use.

  \item[\textbf{Yellow-OCR}]
  Color-filtered OCR baseline. Isolates yellow pixels via HSV masking before
  running easyocr. Tests whether color-aware preprocessing can substitute for
  reasoning or executable recovery.

  \item[\textbf{MCP Skill}]
  Heuristic agentic baseline (\texttt{solve\_captcha.py}): regex-based
  instruction parsing (S1), \texttt{ffmpeg} execution (S2), and a multi-method
  extraction chain (S3).

  \item[\textbf{Claude Haiku / Sonnet / Opus}]
  API-first frontier multimodal systems used to probe the gap between prepared
  extraction and end-to-end recovery.

  \item[\textbf{GPT-5.4}]
  Tool-using frontier model evaluated through the same full-pipeline protocol.
\end{description}

\subsection{Current results}

\begin{table}[h]
\centering\small
\caption{S3 key extraction from ground-truth corrected frames ($n{=}5$ per
level). Best per column in \textbf{bold}.}
\label{tab:s3}
\begin{tabular}{lrrrr|r}
\toprule
Model & L1 & L2 & L3 & L4 & Total \\
\midrule
OCR-only       & 20\% & 25\% & 30\% & 25\% & 25\% \\
Yellow-OCR     & 33\% & --   & --   & --   & --   \\
MCP Skill      & 20\% & 25\% & 30\% & 25\% & 25\% \\
Haiku          & \textbf{100\%} & 60\% & 40\% & \textbf{100\%} & 75\% \\
Sonnet         & \textbf{100\%} & \textbf{100\%} & \textbf{100\%} & \textbf{100\%} & \textbf{100\%} \\
Opus           & \textbf{100\%} & \textbf{100\%} & \textbf{100\%} & \textbf{100\%} & \textbf{100\%} \\
\bottomrule
\end{tabular}
\end{table}

\begin{table}[h]
\centering\small
\caption{Full pipeline results (video $\to$ key). This table is intentionally
left incomplete in the current draft; filling it is a required next experiment
because the benchmark's main claim concerns the gap between S3-only extraction
and end-to-end recovery.}
\label{tab:main}
\begin{tabular}{lrrrrrr}
\toprule
Model & L1 EM & L2 EM & L3 EM & L4 EM & EM (all) & PipeF1 \\
\midrule
MCP Skill     & -- & -- & -- & -- & -- & -- \\
Haiku         & -- & -- & -- & -- & -- & -- \\
Sonnet        & -- & -- & -- & -- & -- & -- \\
Opus          & -- & -- & -- & -- & -- & -- \\
GPT-5.4       & -- & -- & -- & -- & -- & -- \\
\bottomrule
\end{tabular}
\end{table}

\paragraph{Overview.}
Table~\ref{tab:s3} reports S3 extraction accuracy in isolation: each model
receives a corrected frame containing the yellow key and must read it. Frontier
VLMs (Opus, Sonnet) achieve perfect accuracy, while Haiku drops on smaller,
randomly positioned keys. OCR baselines perform poorly due to background text
interference and clutter.

Crucially, S3 in isolation is \emph{not} the binding bottleneck for frontier
models. The benchmark's real difficulty lies in the full recovery policy:
agents must infer the inverse transformation program, execute it correctly, and
identify the correct temporal window containing the key. The missing full
pipeline table is therefore not a cosmetic omission; it is the central
remaining experiment required to validate the benchmark's main claim.

\subsection{Video-temporal diagnostic results}

We also ran a controlled DAVIS-2017 video-temporal split in which the model
must aggregate two consecutive frames and invert a symbolic fragment
transformation. This split is not yet the full inverse-\texttt{ffmpeg}
benchmark; it is the current fast iteration domain for studying temporal
aggregation, solvability, and dense process metrics. The current pool contains
78 generated tasks with process metrics and 233 model responses across three
CLI backends. All modern samples in this slice use dynamic DAVIS frame
sequences and pass the current ``likely-readable'' solvability gate.

\begin{table}[h]
\centering\small
\caption{Current DAVIS video-temporal diagnostic slice. EM is strict exact
match. Missing counts empty/NONE final answers. Partial is dense process
partial credit. Values are means over response rows; $\pm$ is population
standard deviation.}
\label{tab:video-temporal-current}
\begin{tabular}{lrrrr}
\toprule
Backend & $n$ & EM $\uparrow$ & Missing $\downarrow$ & Partial $\uparrow$ \\
\midrule
Claude CLI Haiku & 78 & $14.1{\pm}34.8$ & $1.3{\pm}11.2$ & $43.1{\pm}28.6$ \\
Codex CLI GPT-5.3 Spark & 77 & $0.0{\pm}0.0$ & $100.0{\pm}0.0$ & $6.5{\pm}5.6$ \\
Gemini CLI Flash Lite & 78 & $0.0{\pm}0.0$ & $98.7{\pm}11.2$ & $5.0{\pm}5.7$ \\
\bottomrule
\end{tabular}
\end{table}

At the task level, 11 of 78 tasks are \emph{strict-good}: at least one model
achieves EM. A larger set of 45 tasks is \emph{semantic-good}, meaning either
EM or maximum process partial credit at least 0.25. Finally, 77 tasks have at
least one non-missing response with partial credit at least 0.15, making them
usable for debugging even when no backend solves them exactly. The gap between
strict-good and semantic-good is the most important observation: the split
contains many failures that are not random, but partially correct trajectories.

The strongest current backend is Claude CLI Haiku, which often identifies the
operation and produces structured intermediate evidence but still loses
characters or over-reads distractors on hard samples. Codex Spark and Gemini
Flash Lite mostly fail by returning no final answer in this setup, although
their traces occasionally receive small operation/tool credit. This suggests
that backend/tool framing and response protocol are part of the benchmark
surface, not merely implementation details.

\subsection{Required experiments for the camera-ready version}

The current draft still lacks the analyses that best express CIPHER's identity
as a diagnostic benchmark:
\begin{itemize}
  \item scaling curves over chain length, instruction ambiguity, distractor
        density, and temporal sparsity;
  \item paired S3-only vs full-pipeline comparisons with an explicit
        execution-gap metric;
  \item stage-attribution plots for S1, S2, and S3 across model families; and
  \item a curated failure-seed suite with traces.
\end{itemize}

The video-temporal slice supplies the first version of this failure-seed suite:
12 seedbank genomes were selected by validity, non-missing rate, process
partial credit, frame-localization credit, transform similarity, final-answer
similarity, and diversity. The selected genomes cover all four fragment
operations and range from easy validity-tuned settings (4--6 frames, low noise,
large bright text) to hard-mined settings (7--8 frames, high jitter, distractors,
and \texttt{swap\_reverse\_chars}).
